\chapter{Hushwater: A Composed Board Adapted Into a Place}
\label{ch:ctw}

\section{The Opposite Experiment}

Whitegape answers the question \emph{what does an author decide when nothing is
given}. Hushwater answers the harder one. Its arrangement arrived already made,
emitted by the composer described in \cref{ch:composer} from four numbers and a
seed, and the interesting question is what work is left once that has happened.

The answer is most of the work, and none of the same work. The composer settles
where the rectangles go and what each is for. It settles nothing about what the
board is of, what its ground does, what a player sees, or what any of it is
made of, because it has no representation for any of those. A composed board is
flat, grey and rectangular, and it is a complete description of an arrangement.

Both boards are drawn side by side in \cref{fig:composed-card}, which this chapter
is the commentary on. The hub's ring, its yard hole, the frontline's single run and
both wool families survive the adaptation unchanged.

\section{Reproducing the Starting Point}

The board begins at a descriptor of six fields, and the descriptor is the whole
of what has to be kept for the starting point to be reproducible:

\begin{lstlisting}[language=json]
GET  /api/compose?players=18&teams=2&symmetry=rot_180&seedStart=5&count=1
POST /api/compose/pin  { players: 18, teams: 2, symmetry: "rot_180",
                         cell: 4, seed: 5, composerVersion: "body-first-1",
                         schema: 1, index: 0 }
\end{lstlisting}

The composer version is in that descriptor for a reason that the previous
chapter's subject also demonstrates: the composer is not frozen, and a change
to how it lays out a board moves the geometry a seed produces. The version
string is what makes \emph{seed 5} a stable address rather than a description
of whatever the composer did on the day. Three of the hundred and twenty boards
the repository records moved when a defect in the plan compiler was fixed
during the writing of this paper, and the version went from \id{body-first-1}
to \id{body-first-2} in the same commit.

Before anything was changed, the pinned card was read back and written to disk.
Every comparison in this chapter is against that file rather than against a
recollection of it.

\subsection{The Player Count Names a Band, Not a Budget}

The author asked the composer for fourteen, sixteen, eighteen, twenty and
twenty-one players and received byte-identical boards. That is not a defect. The
count selects a size band, and micro spans fourteen to twenty-one, so every
request inside the band composes the same board at the same seed.

Eighteen was kept because it is the middle of the band and carries no other
meaning. What the number does carry is into the export: \id{maxPlayers} becomes
every team's maximum, so the board ships as eighteen against eighteen rather
than nine against nine, and its 8864 cells work out at about 246 square blocks
a player against the band's measured 242.

\section{What Survived the Adaptation}

The most useful way to read an adaptation is by what it did \emph{not} touch,
because those are the places where the composer's answer was accepted.

The hub kept its ring body and both of its yard holes, sixteen cells each, at
their composed narrowest crossing of sixteen blocks. The frontline kept its
single form and its rectangles: \route{POST /plan/inspect} reads
\id{widthBlocks 32, profile straight} on both teams after the adaptation, as it
did before. Both wool approaches kept the \id{i} family they were emitted with.

The ring in particular is the thing the author records as unable to have been
invented as convincingly. It comes with a sixteen-by-sixteen enclosed yard,
which is a hole players go round, and going round a hole is the rotation device
the composer's own law names in \id{CT8}. An author drawing rectangles from
nothing tends not to produce an enclosure, because an enclosure looks like a
mistake until it is played.

\begin{table}[htbp]
\caption{Hushwater as composed against Hushwater as shipped. Every composed
figure is off the pinned card; every adapted figure is a read off the built
board.}
\label{tab:hushwater-delta}
\small
\begin{tabularx}{\textwidth}{L{4.0cm} L{3.6cm} Y}
\toprule
& \textbf{Composed} & \textbf{Adapted} \\
\midrule
Evaluator score & 0.000, no hard term & 0, \id{valid: true} \\
Plan lint & 1 complaint (\id{WL12}) & 1 complaint (\id{EL1}) \\
Unit land & 250 cells against a 226.4 budget & 265 cells \\
Mid land & 0 cells against a 50.3 budget & 24 cells \\
Hub & 1 box, 128 land / 144 footprint & 116 land, three arms re-cut \\
Frontline & 1 box, 64 / 72 & unchanged \\
Wool & 2 boxes, 42 / 48 & 2 boxes, 57 land \\
Spawn & 1 box, 16 / 16 & 16 cells, moved 16 blocks west \\
Pieces, zones, walls & 12, 1, 0 & 17, 3, 1 \\
Distinct piece surfaces & 1 & 7 \\
Strait & 32 blocks, one hop, empty & 40 blocks with a bank in it \\
\bottomrule
\end{tabularx}
\end{table}

\section{The One Row That Is the Whole Argument}

The last row of \cref{tab:hushwater-delta} is the one to read twice. The
composed board has \textbf{one} distinct piece surface. The shipped board has
\textbf{seven}.

A composed board is flat because the composer has no height model, no theme and
no material. Its \id{globals.surface} is 9 and not one piece overrides it. The
adaptation gives the board a ladder from the gill at the bottom to the mine head
at the top, and seven surfaces is what that ladder is, stated at the level where
the studio will enforce it.

Everything an observer would call the character of the map lives in the
difference between those two numbers. The composer decided where the wool rooms
are and how far apart; it had nothing to say about whether the ground between
them falls forty feet.

\section{What the Adaptation Changed, and Why Each Change Was Forced}

Five changes are substantive. Each was made for a stated reason, and in four of
the five the reason is a rule band rather than a preference.

\subsection{The Unit Moved One Cell Back}

Every piece was pushed one cell in $z$, taking the strait from 32 blocks to 40.

The reason is that the author wanted a bank in the gill, and a mid island splits
one crossing into two hops. \id{G5} wants each hop between 10 and 20 blocks. On
a 32-block strait an eight-block island leaves two twelve-block hops but is too
thin to fight on. At 40 the island can be eight deep in two staggered legs and
still leave twelve and twenty.

The strait itself is the composer's idea and the author moved its ends rather
than its idea. This is the only change in the run that altered composed
geometry, and the journal records it as such.

\subsection{The Crossing Acquired a Bank}

The composed mid is empty: zero cells against a budget of 50.3. The adaptation
puts a Z of two neutral legs in it, each the other's \id{rot\_180} image,
meeting at the origin.

What that does to the crossing is the interesting part. The composed board
offers one hop of 32 blocks. The adapted board offers two, of twelve and twenty,
and \id{rot\_180} hands each team the short one on the opposite hand. A single
symmetric crossing becomes an asymmetric pair that is nonetheless perfectly
fair, because the asymmetry is the same for both teams read from their own end.

\subsection{Two Rooms Acquired Aprons, and One Acquired a Wall}

The two wool approaches keep their \id{i} family and are given opposite
characters: the west one open across moor, the east one a lane behind a bedrock
wall five courses above it, on sixteen blocks of interface.

The wall is the board's only defence wall and the composed board has none. It is
also the reason the two rooms are not the same room twice, which a symmetric
board otherwise tends to produce: both teams get one open approach and one
covered one, and which is which depends on the end they start from.

\subsection{A Bing Was Cut Off the Team's Own Ground}

A bing is a spoil heap. Hushwater's is a pad cut off the team's own land and
bridged back to it, and it is there because a board wants somewhere a player
has a reason to go that is not on the way to anything else.

The flow read confirms exactly that and is worth quoting in its own terms: 368
of 8864 blocks, four per cent, sit off every land route, and 144 of those are
the bing. A pad no land route passes is what a team transient link is, so the
figure is the design working rather than a fault.

That number is also the right comparison to make with Whitegape, where the same
read returned just under a fifth of the board. Hushwater's dead ground is one
deliberate pad; Whitegape's is three unvisited corners of fell.

\subsection{The Spawn Moved Sixteen Blocks West}

The composer's spawn box is sixteen cells and stays sixteen cells. It moves onto
the board's centre line, which is a decision about what a player sees on
spawning rather than about area, and the composer has no term for that.

\section{Where the Studio Said No}

Ten distinct rule ids refused or declined something during this run.
\Cref{tab:hushwater-refusals} lists each with the author's response, and the
pattern in the right-hand column is the substance of this section: a refusal is
rarely answered by arguing with the rule.

\begin{table}[htbp]
\caption{Every rule that stopped or dropped something on the Hushwater run.}
\label{tab:hushwater-refusals}
\small
\begin{tabularx}{\textwidth}{L{1.9cm} L{5.4cm} Y}
\toprule
\textbf{Id} & \textbf{What it said} & \textbf{What was done} \\
\midrule
\id{G2} & zone \id{staithe-link-e} corridor width 8 < 10, hard, score 2000 &
  every build zone widened to 12 blocks or more \\
\id{G5} & gap hop 8 outside 10..20, hard &
  the pad moved so both hops are 12 to 16 \\
\id{ST2} & the iron cube does not stand inside a spawn piece &
  moved onto the spawn-role piece \\
\id{WX8} & the cube needs 3$\times$3 inside the piece and 2 clear of the shell &
  the iron dropped: \id{WX1} insets the shell one block from the piece, so the
  three cannot hold on a plan-compiled spawn \\
\id{DR-DOC} & field \id{pave.jitter} could not be read &
  \id{jitter} is an integer percentage, not a fraction; fixed everywhere,
  including inside the themes \\
\id{PT4} & a material sampled in the plane reads as vertical stripes &
  a \id{rise} given to the stair and lane materials \\
\id{HS4} & beams are dark oak and the post is spruce &
  beams cut from the same wood as the posts \\
\id{SK14} & two override adds state a top the relief solves through &
  both given \id{relief\_scope: "exclude"} \\
\id{SK13} & \id{spill-yard} draws nothing over 130 columns; \id{void-1-cut}
  takes them away &
  the brush moved off the composer's yard hole. A hole is never scenery \\
\id{DR} family & seven declines on the first dressing pass &
  each seat moved by reading \route{sketch/seats} \\
\bottomrule
\end{tabularx}
\end{table}

Two complaints ship standing, and both are answered by a read rather than by a
change. \id{EL1} reports that \id{frontline-t1} steps three blocks to
\id{hub-t2}, which is true of the plan and false of the world: the plan tier
walks pieces flat and cannot see the relief, and the built crossing transects at
\emph{rises 6, falls 1, worst step 1, no barrier, no scramble, walked end to
end}. \id{SK27} reports seven plateaus carrying three paints; two of those
boundaries are what the board is about, and the third is crossed deliberately by
four brushes drawn over the risers rather than inside a plateau.

\section{The Defect This Board Found}

Hushwater is also the board that exposed a studio defect.
\Cref{ch:refusals} sets out the code structure that produced it; this section is
the sequence by which it was found, which is the part that generalises to any
board.

The map author, looking at the built world, observed that one wool room's
entrance redstone appeared on its front face only. \id{ST1} says the entrance
line lies on the last block row at \emph{each} of the room's entry interfaces,
and \route{POST /plan/inspect} reported two for that room: an eight-block flush
seam to \id{wool-a-t1}, classified \id{narrow}, and a twelve-block one to
\id{wool-a-apron}, classified \id{land}. Its mirror image carried two lines. This
one carried one.

\id{ContactGraph} defines the predicate exactly once, and it answers \id{Land}
\emph{or} \id{Narrow}: an interface under the ten-block corridor minimum is a
doorway rather than a wall. The \id{woolRoom} flag each segment carries is
raised through that predicate. Two sites in the plan compiler re-derived the
test inline as \id{Kind: ContactKind.Land}, so a segment was flagged by one
reading and then discarded by another that disagreed with the flag it was
reading.

The consequence was larger than a missing marker. The second of those two sites
is what the exporter cuts cage doors from, so a room whose only interface is
narrower than ten blocks compiled with no door at all, in a room the validator
considers reachable. Nothing refused it and nothing warned.

\begin{ruling}{What the author saw, and what it turned out to be}
``The redstone line is sort of not appearing all the way around it. It only
seems to appear at the front face and not like the other sides, which I'm not
sure why that's happening.'' The observation was of one marker. The defect was a
room that could not be entered.
\end{ruling}

\section{The Division of Labour on This Board}

\Cref{tab:hushwater-labour} is the counterpart of the table that closes
\cref{ch:dtm}. The rows are the same rows; the middle column has grown.

\begin{table}[htbp]
\caption{Hushwater, stage by stage. The composer's column exists only on this
board.}
\label{tab:hushwater-labour}
\small
\begin{tabularx}{\textwidth}{L{2.2cm} L{4.2cm} L{4.2cm} Y}
\toprule
\textbf{Stage} & \textbf{The composer} & \textbf{\studio{The studio}} &
\textbf{\agent{The author}} \\
\midrule
Arrangement & hub form, frontline form, both wool families, the land budget,
  the 32-block strait & \id{G2}, \id{G5}, the evaluator's gate &
  which of 48 boards, and one cell of push \\
Elevation & nothing & the solved field & seven surfaces and every mark \\
Paint & nothing & which block on which cell & three themes and where the bands
  cut \\
Structures & nothing & storey heights, roof pitch, door side & the style, and
  what stands where \\
Dressing & nothing & the legal seats, ten declines & which seats to take \\
Objectives & where the wools are & the region and the filter & what the board is
  played for \\
\bottomrule
\end{tabularx}
\end{table}
